\documentclass{amsart}
\usepackage{amsmath}
\usepackage{amsfonts}
\usepackage{amssymb}
\usepackage{graphicx}
\usepackage[miktex]{gnuplottex}
\usepackage{epstopdf}
\usepackage{minted}
\usepackage{color,soul}
\usepackage{listings}
\usepackage{booktabs}
\usepackage{hyperref}
\newtheorem{theorem}{Theorem}
\theoremstyle{plain}
\newtheorem{acknowledgement}{Acknowledgement}
\newtheorem{algorithm}{Algorithm}
\newtheorem{axiom}{Axiom}
\newtheorem{case}{Case}
\newtheorem{claim}{Claim}
\newtheorem{conclusion}{Conclusion}
\newtheorem{condition}{Condition}
\newtheorem{conjecture}{Conjecture}
\newtheorem{corollary}{Corollary}
\newtheorem{criterion}{Criterion}
\newtheorem{definition}{Definition}
\newtheorem{example}{Example}
\newtheorem{exercise}{Exercise}
\newtheorem{lemma}{Lemma}
\newtheorem{notation}{Notation}
\newtheorem{problem}{Problem}
\newtheorem{proposition}{Proposition}
\newtheorem{remark}{Remark}
\newtheorem{solution}{Solution}
\newtheorem{summary}{Summary}
\numberwithin{equation}{section}
%--------------------------------------------------------

\begin{document}

\title{SVI model in ORE}

\date{3 February 2026}

\begin{abstract}
Description of SVI model implementation in ORE.
\end{abstract}
\maketitle

% \tableofcontents

\section{SVI Model}\label{sec:model}

\subsection{Specification}

The SVI model is given by (see e.g. \cite{Gatheral_2004})

\begin{equation}
w(k; a,b,\rho,m,\sigma) = a + b \left( \rho (k - m) + \sqrt{(k - m)^2 + \sigma^2} \right)
\end{equation}

with

\begin{itemize}
\item $w$ the total implied variance, i.e. $w = \sigma_{imp}^2 T$ with implied volatility $\sigma_{imp}$ and
  option time to maturity $T$
\item $k$ the log-moneyness, i.e. $k = \ln(K/F)$ with strike $K$ and forward $F$
\item $a$ controls the overall level of the variance curve
\item $b \ge 0$ controls the slope of the variance curve
\item $\rho \in [-1,1]$ controls the skewness of the variance curve
\item $m$ controls the horizontal shift of the variance curve
\item $\sigma > 0$ controls the curvature of the variance curve
\end{itemize}

\subsubsection{Raw SVI}\label{sec:rawsvi}

The above is the so-called ``raw'' SVI formulation. In the
following we will refer to the $5$ parameters $(a,b,\rho,m,\sigma)$ as the ``raw SVI parameters''.
There are two other SVI variants known from the literature.

\subsubsection{Natural SVI}\label{sec:naturalsvi}

The ``natural'' SVI formulation uses parameters $(\Delta, \mu, \rho, \omega, \zeta)$ related to the raw parameters via
\begin{eqnarray}
a &=& \Delta + \frac{\omega}{2} (1 - \rho^2) \\
b &=& \frac{\zeta \omega}{2} \\
m &=& \mu - \frac{\rho \omega}{\zeta} \\
\sigma &=& \frac{\omega}{\zeta}
\end{eqnarray}

The total variance is therefore given by
\begin{equation}
w(k; \Delta, \mu, \rho, \omega, \zeta) = \Delta + \frac{\omega}{2} \left( 1 + \zeta \rho (k - \mu) + \sqrt{(\zeta (k - \mu) + \rho)^2 + 1 - \rho^2} \right)
\end{equation}

\subsubsection{SVI Jump-Wings}\label{sec:svijw}

The ``SVI-JW'' formulation uses parameters $(v_t, \psi_t, p_t, c_t, \tilde{v}_t)$ related to the raw parameters via

\begin{eqnarray}
v_t &=& \frac{a + b (- \rho m + \sqrt{m^2 + \sigma^2})}{t} \\
\psi_t &=& \frac{1}{\sqrt{w_t}} \frac{b}{2} (-\frac{m}{\sqrt{m^2 + \sigma^2}} + \rho)  \\
p_t &=& \frac{1}{\sqrt{w_t}}  b (1 - \rho) \\
c_t &=& \frac{1}{\sqrt{w_t}}  b (1 + \rho) \\
\tilde{v}_t &=& \frac{1}{t}(a + b \sigma \sqrt{1 - \rho^2})
\end{eqnarray}

Closed-form translation from SVI-JW parameters to raw SVI parameters exists only under certain conditions, see
\cite{Gatheral_2004} for details.

\subsubsection{Surface SVI}\cite{Gatheral_2012}\label{sec:ssvi}

The Surface SVI model is an extension of the natural parameterization, with time-dependent parameters
$(\theta_t, \varphi(\theta_t))$ and non-time-dependent parameters $\rho$, via

\begin{equation}
w(k,t) = \frac{\theta_t}{2} \left( 1 + \varphi(\theta_t) \rho k + \sqrt{(\varphi(\theta_t) k + \rho)^2 + 1 - \rho^2} \right)
\end{equation}

with

\begin{eqnarray}
\theta_t &=& \sigma_{BS}^2(0,t) t \\
\varphi: \mathbb{R}^+ &\rightarrow& \mathbb{R}^+ \text{ (smooth function)}
\end{eqnarray}

Sufficient conditions for absence of calendar and butterfly arbitrage are, $\forall t, \theta > 0$,

\begin{enumerate}
\item $\partial_t \theta_t \ge 0$
\item $0 \le \partial_\theta (\theta \varphi(\theta)) \le \frac{1}{\rho^2} (1 + \sqrt{1 - \rho^2})$
\item $\theta \varphi(\theta) (1 + |\rho|) < 4$
\item $\theta \varphi(\theta)^2 (1 + |\rho|) < 4$
\end{enumerate}

\subsubsection{Surface SVI - Heston-like Parameterization}\label{sec:ssvi_heston}
A popular choice for $\varphi$ is the Heston-like parameterization
\begin{equation}
\varphi(\theta) = \frac{1}{\lambda \theta}(1 - \frac{1 - \exp(-\lambda \theta)}{\lambda \theta})  \end{equation}
with $\lambda > 0$ a constant parameter.

The model guarantees absense of calendar arbitrage, and sufficient condition for absense of butterfly arbitrage is
\begin{equation}
\lambda \ge \frac{1+|\rho|}{4}
\end{equation}

\subsubsection{Surface SVI - Power-law Parameterization}\label{sec:ssvi_power_law}

Another choice for $\varphi$ is the power-law parameterization
\begin{equation}
\varphi(\theta) = \eta \theta^{-\gamma}
\end{equation}
with $\eta > 0$ and $0 < \gamma < 1$ constant parameters

The model guarantees absense of calendar arbitrage, and sufficient conditions for absense of butterfly arbitrage are
\begin{eqnarray}
\theta &\le& \left(\frac{4}{\eta (1 + |\rho|)}\right)^{1-\gamma}  \\
\theta &\le& \left(\frac{4}{\eta^2 (1 + |\rho|)}\right)^{\frac{1}{1-2\gamma}} \quad \text{for } 0<\gamma<0.5 \\
\theta &\ge& \left(\frac{4}{\eta^2 (1 + |\rho|)}\right)^{\frac{1}{1-2\gamma}} \quad \text{for } 0.5<\gamma<1 \\
4 &\ge& \eta^2 (1 + |\rho|) \quad \text{for } \gamma=0.5
\end{eqnarray}

\subsubsection{Extended Surface SVI}\cite{Hendriks_2017}\label{sec:essvi}

The extended Surface SVI model is an extension of the Surface SVI model, replacing the static parameter $\rho$ with a time-dependent parameter
$\rho_t$, via
\begin{equation}
w(k,t) = \frac{\theta_t}{2} \left( 1 + \varphi(\theta_t) \rho_t k + \sqrt{(\varphi(\theta_t) k + \rho_t)^2 + 1 - \rho_t^2} \right)
\end{equation}

\subsection{Model Variants}

ORE supports the following SVI model variants:

\begin{itemize}
\item Gatheral2004SviRaw: SVI with Raw parameterization (see \ref{sec:rawsvi})
\item Gatheral2004SviNatural: SVI with Natural parameterization (see \ref{sec:naturalsvi})
\item Gatheral2004SviJw: SVI with Jump-Wings parameterization (see \ref{sec:svijw})
\item Gatheral2012SsviHeston: Surface SVI with Heston-like parameterization (see \ref{sec:ssvi_heston}), arbitrage-free
\item Gatheral2012SsviPowerLaw: Surface SVI with power-law parameterization (see \ref{sec:ssvi_power_law}), arbitrage-free
\item HendriksMartini2017EssviFirstPowerLaw: Extended Surface SVI with first power-law parameterization (see \ref{sec:essvi}), arbitrage-free
\item HendriksMartini2017EssviSecondPowerLaw: Extended Surface SVI with second power-law parameterization (see \ref{sec:essvi}), arbitrage-free
\item CorbettaEtAl2019Essvi: Extended Surface SVI (see \ref{sec:essvi}) using robust calibration \cite{Corbetta_2019}, arbitrage-free,
\item Mingone2022EssviGJ: Extended Surface SVI (see \ref{sec:essvi}) with global calibration and Gatheral Jim condition \cite{Mingone_2022}, arbitrage-free,
\item Mingone2022EssviMM: Extended Surface SVI (see \ref{sec:essvi}) with global calibration and Martini Mingone condition \cite{Mingone_2022}, arbitrage-free,
\end{itemize}

\begin{remark}
The arbitrage-free claims above apply to equity, FX, and commodity options, where the pricing numeraire is the
risk-free discount factor $B(0,T)$.  For CDS options the correct numeraire is the risky annuity $\text{rpv01}(T)$,
which is default-contingent and strictly smaller than the risk-free annuity.  Both the calendar-spread and butterfly
arbitrage conditions are derived in the Black-76 / risk-free-discount setting and do not carry over exactly to CDS
surfaces: even if total variance is non-decreasing in $t$, the declining $\text{rpv01}(T)$ can render a calendar spread
negative; similarly, butterfly-arbitrage-free densities under the $T$-forward measure need not be non-negative under the
annuity measure.  Both conditions are approximately satisfied for investment-grade names where
$\text{rpv01}(T) \approx B(0,T) \cdot \text{annuity}$, but become weaker guarantees for high-yield or distressed
credits.
\end{remark}

\subsection{Calibration}

The calibration of the SVI parameters is performed on the market smile points, which are given by

\begin{itemize}
\item option expiry and swap underlying length in the case of credit default swaptions
\item option expiry in the case of equity options, FX options and commodity options 
\end{itemize}

If no free parameters are given, we use the given initial SVI parameters unchanged for the given market point.

If at least one free parameter is given, we run an optimization using the LevenbergMarquardt optimizer on the following
target function $f$.

\begin{equation}
T(x) = \sum_{i=1}^n m_i - s(x, K_i)
\end{equation}

with

\begin{itemize}
\item $x$ the input SVI parameters
\item $m_i$ the $ith$ market quote, converted to the preferred output type of the resolution method
\item $s(x,K)$ the output of the resolution method for parameters $x$ and strike $K$
\end{itemize}

The optimization is run from user-provided initial SVI parameters. If the target function $T$ on this first run is
below a user-provided ``early exit'' error threshold $t_A$, i.e.

$$
T(x) < t_A
$$

the result of the optimization will be accepted. Otherwise more optimization runs will be performed using random initial
guesses generated from a Halton sequence.

If for any of the subsequent optimization runs, the error lies below $t_A$ the corresponding minimum will be
accepted. If for a given maximum number of runs $N$ the error never lies below the given threshold, the best found
solution with target value $T(x_b)$ is compared against the maximum acceptable threshold $t_B$ and if

$$
T(x_b) < t_B
$$

the solution $x_b$ is accepted. Otherwise the calibration is marked as ``unsuccessful''. After all market points have been
calibrated, we have either a 1D arrays or a 2D matrices of calibrated SVI
parameters with possibly missing entries from unsuccessful calibrations. As a last step, these missing values are
interpolated linearly (1D case) or using Laplace Interpolation in 2D where the x and y axis are indexed by the option
time to maturity and underlying swap length (i.e. the interpolation grid is not equidistant in general). The
interpolated values are updated if not admissible for the SVI model.

\subsection{Calibration for SVI}

Applicable to model variants Gatheral2004SviRaw, Gatheral2004SviNatural, Gatheral2004SviJw.

For each market smile point the free SVI parameters are calibrated to quotes of the given market smile. If the market
smile quotes are not given in the preferred output type of the resolution method, we convert the market quote to the
preferred output type first.

\subsubsection{Interpolation}

SVI parameters for option expiry times / underlying swap lengths that do not lie on the market quotation grid will be
interpolated linearly (1D case) or bilinearly (2D case) with flat extrapolation.

\subsubsection{Extrapolation} Flat extrapolation is used for option expiry times / underlying swap lengths that lie outside the market quotation grid.

\subsection{Calibration for Surface SVI}

Applicable to model variants Gatheral2012SsviHeston, Gatheral2012SsviPowerLaw.

Global calibration is applied to the parameters $\theta_i$'s, $\rho$, $\lambda$ (for Gatheral2012SsviHeston), $\eta$ and $\gamma$ (for Gatheral2012SsviPowerLaw),
with constraints to guarantee absence of calendar and butterfly arbitrage.

\subsubsection{Interpolation}

Price of options with expiry times that do not lie on the market quotation grid will be interpolated in total variance space,

\begin{equation}
\frac{C_t}{K_t} = \alpha_t \frac{C_i}{K_i} + (1 - \alpha_t) \frac{C_{i+1}}{K_{i+1}}
\end{equation}

where

\begin{equation}
\alpha_t = \frac{\sqrt{\theta_{t_{i+1}}} - \sqrt{\theta_t}}{\sqrt{\theta_{t_{i+1}}} - \sqrt{\theta_{t_i}}}
\end{equation}

with $\theta_t$ linearly interpolated between $\theta_{t_i}$ and $\theta_{t_{i+1}}$ in $t$.

\begin{remark}
The prices $C_i$ and $C_{i+1}$ are computed as forward-premium (discount $= 1$) Black
formulas.  For CDS options the engine discounts with the risky annuity $\text{rpv01}(T)$
rather than a risk-free factor, so the no-arbitrage guarantee is approximate; for
investment-grade names the difference is small.
\end{remark}

\subsubsection{Extrapolation}
For $t < t_1$, we re-use the equation for interpolation with

$(C_i,K_i)$ being the intrinsic value and strike of the option and $\theta_i$ being zero.
$(C_{i+1},K_{i+1})$ being the price and strike of the option with expiry $t_1$ and $\theta_{i+1}$ being the total variance at $t_1$

For $t > t_N$, we extrapolate in total variance space with
\begin{equation}
w(k,\theta_t) = w(k,\theta_{t_N}) + \theta_t - \theta_{t_N}
\end{equation}

\begin{remark}
The early-extrapolation price $C_1$ is a forward-premium (discount $= 1$) Black formula;
the same CDS caveat as in the interpolation section above applies.
Late extrapolation ($t > t_N$) works entirely in total variance space and does not involve
a Black formula, so it is unaffected by the choice of discount numeraire.
\end{remark}

\subsection{Calibration for Extended Surface SVI}

Applicable to model variants HendriksMartini2017EssviFirstPowerLaw, HendriksMartini2017EssviSecondPowerLaw.

\subsubsection{Interpolation}
Work in progress
\subsubsection{Extrapolation}
Work in progress

\subsection{Robust calibration for Extended Surface SVI}

Applicable to model variant CorbettaEtAl2019Essvi.

\cite{Corbetta_2019} proposes a robust calibration method for a re-parametrized Extended Surface SVI model, with new model parameters
$k^*, \theta^*, \rho, \psi \in \mathbb{R} \times \mathbb{R}^+ \times (-1, 1) \times (0, \infty)$
$\theta = \theta^* - \rho \psi k^*$
where the data point $(k^*, \theta^*)$ represents the log-forward moneyness and total implied variance closest to the ATM (forward).

Butterfly arbitrage is avoided by imposing the following constraints on the model parameters:
$0 < \psi \le \frac{-2\rho k^*}{1 + |\rho|} + \sqrt{\frac{4 \rho^2 (k^*)^2}{(1 + |\rho|)^2} + \frac{4 \theta^*}{1 + |\rho|}}$

Calendar arbitrage is avoided by imposing the following constraints on the model parameters:
$\theta > \underline{\theta} \quad \text{and} \quad \psi > \underline{\psi}$

where $\underline{\theta}$ and $\underline{\psi}$ are the calibrated $\theta$ and $\psi$ parameters associated to the previous slice

\subsubsection{Interpolation}

For two slices with parameters $\theta_1, \psi_1, \rho_1$ and $\theta_2, \psi_2, \rho_2$ the interpolated parameters for a slice with time to maturity $t$ between the two slices are given by

\begin{eqnarray}
\theta_t &=& (1 - \lambda) \theta_1 + \lambda \theta_2 \\
\psi_t &=& (1 - \lambda) \psi_1 + \lambda \psi_2 \\
\psi_t \rho_t &=& (1 - \lambda) \psi_1 \rho_1 + \lambda \psi_2 \rho_2
\end{eqnarray}

where $\lambda = \frac{t - T_1}{T_2 - T_1}$ is the interpolation weight.

\subsubsection{Extrapolation}

For $t < T_1$,
\begin{eqnarray}
\theta_t &=& \lambda \theta_1 \\
\psi_t &=& \lambda \psi_1 \\
\rho_t &=& \rho_1
\end{eqnarray}

where $\lambda = \frac{t}{T_1}$ is the extrapolation weight.

For $t > T_N$,

\begin{eqnarray}
\theta_t &=& \theta_N + u(t) \\
\psi_t &=& \psi_N \\
\rho_t &=& \rho_N
\end{eqnarray}

where $u(t)$ is any continuous increasing function on $[T_N, \infty)$ with $u(T_N) = 0$.

In our implementation we use $u(t) = \frac{\theta_N - \theta_{N-1}}{T_N - T_{N-1}} (t - T_N)$


\subsection{Global calibration for Extended Surface SVI}

Applicable to model variant Mingone2022EssviGJ and Mingone2022EssviMM.

ORE implements the global calibration method for the Extended Surface SVI model proposed by Mingone \cite{Mingone_2022}.
The method guarantees absence of both calendar and butterfly arbitrage across the whole surface by re-parametrization of the model,
with new model parameters

\begin{equation}
\rho_1, ..., \rho_N, \theta_1, a_2, ..., a_N, c_1, ..., c_N \in (-1, 1)^N \times (0, \infty)^N \times (0, 1)^N
\end{equation}

via

\begin{eqnarray}
a_i &=& \theta_i - \theta_{i-1} p_i \\
c_i &=& \frac{\psi_i - A_{\psi_i}}{C_{\psi_i} - A_{\psi_i}} \\
p_i &=& \max \left( \frac{1 + \rho_{i-1}}{1 + \rho_i}, \frac{1 - \rho_{i-1}}{1 - \rho_i} \right) \quad \text{for } i > 1 \\
A_{\psi_1} &=& 0 \\
A_{\psi_i} &=& \psi_{i-1} p_i \quad \text{for } i \ge 1 \\
C_{\psi_1} &=& \min \left( f_1, \frac{f_2}{p_2}, ..., \frac{f_N}{\Pi^N_{j=2}p_j} \right) \\
C_{\psi_i} &=& \min \left( \frac{\psi_{i-1}}{\theta_{i-1}},f_i, \frac{f_{i+1}}{p_{i+1}}, ..., \frac{f_N}{\Pi^N_{j=i+1}p_j} \right) \quad \text{for } i \ge 1
\end{eqnarray}

\begin{eqnarray}
f_i &=& \min \left( \frac{4}{1+|\rho_i|}, \sqrt{f_{GJ}(\theta_i, |\rho_i|)}\right) \quad \text{for } i \ge 1 \text{ or }\\
f_i &=& \min \left( \frac{4}{1+|\rho_i|}, \sqrt{f_{MM}(\theta_i, |\rho_i|)}\right) \quad \text{for } i \ge 1 \\
\end{eqnarray}

where $f_{GJ}$ and $f_{MM}$ are the Gatheral Jim and Martini Mingone functions defined in \cite{Mingone_2022}.

The natural SVI parameters can be recovered via
\begin{eqnarray}
\theta_2 &=& \theta_1 p_2 + a_2, ..., \theta_N = \theta_{N-1} p_N + a_N \\
\psi_1 &=& c_1 (C_{\psi_1} - A_{\psi_1}) + A_{\psi_1}, ..., \psi_N = c_N (C_{\psi_N} - A_{\psi_N}) + A_{\psi_N} \\
\phi_i &=& \frac{\psi_i}{\theta_i}, \quad \rho_i \text{ given by the model parameters}
\end{eqnarray}

Note that logit transformation is applied to $c_i$ internally to map the optimization domain from $(0,1)$ to
$(-\infty, \infty)$.

To maintain the absence of arbitrage across the whole surface, interpolation and extrapolation of the model parameters is performed on natural SVI parameters $\theta_i, \psi_i, \psi_i \rho_i$ for $i = 1, ..., N$.

\subsubsection{Interpolation}

For two slices with parameters $\theta_1, \psi_1, \psi_1 \rho_1$ and $\theta_2, \psi_2, \psi_2 \rho_2$ the interpolated parameters for a slice with time to maturity $t$ between the two slices are given by

\begin{eqnarray}
\theta_t &=& (1 - \lambda) \theta_1 + \lambda \theta_2 \\
\psi_t &=& (1 - \lambda) \psi_1 + \lambda \psi_2 \\
\psi_t \rho_t &=& (1 - \lambda) \psi_1 \rho_1 + \lambda \psi_2 \rho_2
\end{eqnarray}

where $\lambda = \frac{t - T_1}{T_2 - T_1}$ is the interpolation weight.

\subsubsection{Extrapolation}

For $t < T_1$,
\begin{eqnarray}
\theta_t &=& \lambda \theta_1 \\
\psi_t &=& \lambda \psi_1 \\
\rho_t &=& \rho_1
\end{eqnarray}

where $\lambda = \frac{t}{T_1}$ is the extrapolation weight.

For $t > T_N$,

\begin{eqnarray}
\theta_t &=& \theta_N  + \frac{\theta_N - \theta_{N-1}}{T_N - T_{N-1}} (t - T_N) \\
\psi_t &=& \psi_N \\
\rho_t &=& \rho_N
\end{eqnarray}

\section{Parametrization in ORE}

For details on the parametrization see also the ORE User Guide. This documentation focuses on the link between the
parametrization and methodology explained in \ref{sec:model}.

The SVI interpolation is specified as

\begin{minted}[fontsize=\footnotesize]{xml}
      <TimeInterpolation>Mingone2022EssviGJ</TimeInterpolation>
      <StrikeInterpolation>Mingone2022EssviGJ</StrikeInterpolation>
\end{minted}

within a \texttt{StrikeSurface}, \texttt{MoneynessSurface}, or
\texttt{DeltaSurface} element, for equity and commodity volatility surfaces. Both fields must be set to the same SVI variant name. The {\em Input} volatility type is given by

\begin{minted}[fontsize=\footnotesize]{xml}
      <VolatilityType>Lognormal</VolatilityType>
\end{minted}

The displacement $d$ applied in the SVI {\em model} formulation is either read from the input market quotes, or can be
overridden with

\begin{minted}[fontsize=\footnotesize]{xml}
      <ModelShift>0.0</ModelShift>
\end{minted}

for equity options, FX options and commodity options (single displacement $d$ for the whole surface).

\subsection{ParametricSmileConfiguration}

The \texttt{ParametricSmileConfiguration} element specifies the initial values and calibration settings for the SVI
parameters. It is specified within the surface element (e.g. \texttt{StrikeSurface}, \texttt{MoneynessSurface}, or
\texttt{DeltaSurface}) alongside the \texttt{TimeInterpolation} and \texttt{StrikeInterpolation} elements (which must both be set to the same SVI variant name).

The initial values for the optimization for each parameter are specified, either as a single value that is used for all
options or individually for the array of option expiries.

For the Gatheral2004SviRaw model variant, the parameters are $a$, $b$, $\rho$, $m$ and $\sigma$:

\begin{minted}[fontsize=\footnotesize]{xml}
      <Parameters>
        <Parameter>
          <Name>a</Name>
          <InitialValue>0.05</InitialValue>
          <Calibration>Implied</Calibration>
        </Parameter>
        <Parameter>
          <Name>b</Name>
          <InitialValue>0.1</InitialValue>
          <Calibration>Fixed</Calibration>
        </Parameter>
        <Parameter>
          <Name>rho</Name>
          <InitialValue>0.0</InitialValue>
          <Calibration>Calibrated</Calibration>
        </Parameter>
        <Parameter>
          <Name>m</Name>
          <InitialValue>0.5</InitialValue>
          <Calibration>Calibrated</Calibration>
        </Parameter>
        <Parameter>
          <Name>sigma</Name>
          <InitialValue>0.1</InitialValue>
          <Calibration>Calibrated</Calibration>
        </Parameter>
      </Parameters>
\end{minted}

For the Gatheral2004SviNatural model variant, the parameters are $\Delta$, $\mu$, $\rho$, $\omega$ and $\zeta$:

\begin{minted}[fontsize=\footnotesize]{xml}
      <Parameters>
        <Parameter>
          <Name>delta</Name>
          <InitialValue>0.05</InitialValue>
          <Calibration>Implied</Calibration>
        </Parameter>
        <Parameter>
          <Name>mu</Name>
          <InitialValue>0.0</InitialValue>
          <Calibration>Calibrated</Calibration>
        </Parameter>
        <Parameter>
          <Name>rho</Name>
          <InitialValue>0.0</InitialValue>
          <Calibration>Calibrated</Calibration>
        </Parameter>
        <Parameter>
          <Name>omega</Name>
          <InitialValue>0.1</InitialValue>
          <Calibration>Calibrated</Calibration>
        </Parameter>
        <Parameter>
          <Name>zeta</Name>
          <InitialValue>0.1</InitialValue>
          <Calibration>Calibrated</Calibration>
        </Parameter>
      </Parameters>
\end{minted}

For the Gatheral2004SviJw model variant, the parameters are $v_t$, $\psi_t$, $p_t$, $c_t$ and $\tilde{v}_t$:

\begin{minted}[fontsize=\footnotesize]{xml}
      <Parameters>
        <Parameter>
          <Name>v</Name>
          <InitialValue>0.04</InitialValue>
          <Calibration>Implied</Calibration>
        </Parameter>
        <Parameter>
          <Name>psi</Name>
          <InitialValue>0.0</InitialValue>
          <Calibration>Calibrated</Calibration>
        </Parameter>
        <Parameter>
          <Name>p</Name>
          <InitialValue>0.1</InitialValue>
          <Calibration>Calibrated</Calibration>
        </Parameter>
        <Parameter>
          <Name>c</Name>
          <InitialValue>0.1</InitialValue>
          <Calibration>Calibrated</Calibration>
        </Parameter>
        <Parameter>
          <Name>vtilde</Name>
          <InitialValue>0.04</InitialValue>
          <Calibration>Calibrated</Calibration>
        </Parameter>
      </Parameters>
\end{minted}

For the Gatheral2012SsviHeston model variant, the parameters are $\theta$, $\rho$ and $\lambda$:

\begin{minted}[fontsize=\footnotesize]{xml}
      <Parameters>
        <Parameter>
          <Name>theta</Name>
          <InitialValue>0.02</InitialValue>
          <Calibration>Implied</Calibration>
        </Parameter>
        <Parameter>
          <Name>rho</Name>
          <InitialValue>0.0</InitialValue>
          <Calibration>Calibrated</Calibration>
        </Parameter>
        <Parameter>
          <Name>lambda</Name>
          <InitialValue>1.0</InitialValue>
          <Calibration>Calibrated</Calibration>
        </Parameter>
      </Parameters>
\end{minted}

For the Gatheral2012SsviPowerLaw model variant, the parameters are $\theta$, $\rho$, $\eta$ and $\gamma$:

\begin{minted}[fontsize=\footnotesize]{xml}
      <Parameters>
        <Parameter>
          <Name>theta</Name>
          <InitialValue>0.02</InitialValue>
          <Calibration>Implied</Calibration>
        </Parameter>
        <Parameter>
          <Name>rho</Name>
          <InitialValue>0.0</InitialValue>
          <Calibration>Calibrated</Calibration>
        </Parameter>
        <Parameter>
          <Name>eta</Name>
          <InitialValue>0.5</InitialValue>
          <Calibration>Calibrated</Calibration>
        </Parameter>
        <Parameter>
          <Name>gamma</Name>
          <InitialValue>0.5</InitialValue>
          <Calibration>Calibrated</Calibration>
        </Parameter>
      </Parameters>
\end{minted}

For the CorbettaEtAl2019Essvi model variant, the parameters are $k^*$, $\theta^*$, $\rho$ and $\psi$:

\begin{minted}[fontsize=\footnotesize]{xml}
      <Parameters>
        <Parameter>
          <Name>k_star</Name>
          <InitialValue>0.0</InitialValue>
          <Calibration>Implied</Calibration>
        </Parameter>
        <Parameter>
          <Name>theta_star</Name>
          <InitialValue>0.0</InitialValue>
          <Calibration>Implied</Calibration>
        </Parameter>
        <Parameter>
          <Name>rho</Name>
          <InitialValue>0.0</InitialValue>
          <Calibration>Calibrated</Calibration>
        </Parameter>
        <Parameter>
          <Name>psi</Name>
          <InitialValue>0.01</InitialValue>
          <Calibration>Calibrated</Calibration>
        </Parameter>
      </Parameters>
\end{minted}

For the Mingone2022EssviGJ and Mingone2022EssviMM model variants, the parameters are $\rho$, $a$ and $c$ (see
\ref{sec:model}):

\begin{minted}[fontsize=\footnotesize]{xml}
      <Parameters>
        <Parameter>
          <Name>rho</Name>
          <InitialValue>0.0</InitialValue>
          <Calibration>Calibrated</Calibration>
        </Parameter>
        <Parameter>
          <Name>a</Name>
          <InitialValue>0.0</InitialValue>
          <Calibration>Calibrated</Calibration>
        </Parameter>
        <Parameter>
          <Name>c</Name>
          <InitialValue>0.5</InitialValue>
          <Calibration>Calibrated</Calibration>
        </Parameter>
      </Parameters>
\end{minted}

The ``Calibration'' field is used to distinguish

\begin{itemize}
\item ``Fixed'': fixed values that are not calibrated, but prescribed by the user
\item ``Implied'': if the values are implied from market data
\item ``Calibrated'': takes part in the optimization
\end{itemize}

It follows a Calibration section with the maximum number of calibration attempts $N$, the early exit error threshold
$t_A$ and the maximum acceptable error threshold $t_B$:

\begin{minted}[fontsize=\footnotesize]{xml}
      <Calibration>
        <MaxCalibrationAttempts>10</MaxCalibrationAttempts>
        <ExitEarlyErrorThreshold>0.005</ExitEarlyErrorThreshold>
        <MaxAcceptableError>0.05</MaxAcceptableError>
      </Calibration>
\end{minted}

\subsection{Equity Volatility Surface Example (Mingone)}

A complete example for an equity volatility surface using the Mingone2022EssviGJ model variant:

\begin{minted}[fontsize=\footnotesize]{xml}
  <EquityVolatility>
    <CurveId>RIC:.SPX</CurveId>
    <CurveDescription>Lognormal option implied vols for SPX
      using ESSVI Mingone GJ</CurveDescription>
    <Currency>USD</Currency>
    <DayCounter>A365</DayCounter>
    <VolatilityConfig>
      <StrikeSurface priority="0">
        <QuoteType>ImpliedVolatility</QuoteType>
        <VolatilityType>Lognormal</VolatilityType>
        <Strikes>*</Strikes>
        <Expiries>*</Expiries>
        <TimeInterpolation>Mingone2022EssviGJ</TimeInterpolation>
        <StrikeInterpolation>Mingone2022EssviGJ</StrikeInterpolation>
        <Extrapolation>true</Extrapolation>
        <TimeExtrapolation>Flat</TimeExtrapolation>
        <StrikeExtrapolation>Flat</StrikeExtrapolation>
        <ParametricSmileConfiguration>
          <Parameters>
            <Parameter>
              <Name>rho</Name>
              <InitialValue>0.0</InitialValue>
              <Calibration>Calibrated</Calibration>
            </Parameter>
            <Parameter>
              <Name>a</Name>
              <InitialValue>0.0</InitialValue>
              <Calibration>Calibrated</Calibration>
            </Parameter>
            <Parameter>
              <Name>c</Name>
              <InitialValue>0.5</InitialValue>
              <Calibration>Calibrated</Calibration>
            </Parameter>
          </Parameters>
          <Calibration>
            <MaxCalibrationAttempts>50</MaxCalibrationAttempts>
            <ExitEarlyErrorThreshold>0.005</ExitEarlyErrorThreshold>
            <MaxAcceptableError>0.5</MaxAcceptableError>
          </Calibration>
        </ParametricSmileConfiguration>
      </StrikeSurface>
    </VolatilityConfig>
  </EquityVolatility>
\end{minted}



\subsection{FX Volatility Surface Example (Mingone)}

A complete example for an FX volatility surface using the Mingone2022EssviGJ model variant:

\begin{minted}[fontsize=\footnotesize]{xml}
  <FXVolatility>
    <CurveId>USDJPY</CurveId>
    <CurveDescription>Lognormal option implied vols for USD/JPY
      using ESSVI Mingone GJ</CurveDescription>
    <Dimension>Smile</Dimension>
    <SmileType>BFRR</SmileType>
    <SmileInterpolation>Mingone2022EssviGJ</SmileInterpolation>
    <TimeInterpolation>V2T</TimeInterpolation>
    <SmileDelta>25</SmileDelta>
    <Conventions>USD-JPY-FXOPTION</Conventions>
    <Expiries>*</Expiries>
    <FXSpotID>FX/USD/JPY</FXSpotID>
    <FXForeignCurveID>Yield/USD/USD-FedFunds</FXForeignCurveID>
    <FXDomesticCurveID>Yield/JPY/JPY-IN-USD</FXDomesticCurveID>
    <Calendar>JoinHolidays(US settlement, Japan)</Calendar>
    <DayCounter>Actual/365 (Fixed)</DayCounter>
    <ParametricSmileConfiguration>
      <Parameters>
        <Parameter>
          <Name>rho</Name>
          <InitialValue>0.0</InitialValue>
          <Calibration>Calibrated</Calibration>
        </Parameter>
        <Parameter>
          <Name>a</Name>
          <InitialValue>0.0</InitialValue>
          <Calibration>Calibrated</Calibration>
        </Parameter>
        <Parameter>
          <Name>c</Name>
          <InitialValue>0.5</InitialValue>
          <Calibration>Calibrated</Calibration>
        </Parameter>
      </Parameters>
      <Calibration>
        <MaxCalibrationAttempts>50</MaxCalibrationAttempts>
        <ExitEarlyErrorThreshold>0.005</ExitEarlyErrorThreshold>
        <MaxAcceptableError>0.5</MaxAcceptableError>
      </Calibration>
    </ParametricSmileConfiguration>
  </FXVolatility>
\end{minted}

\subsection{Commodity Volatility Surface Example (Corbetta Robust)}

SVI is supported for commodity volatility surfaces using all three 2D surface config types:
\texttt{StrikeSurface} (strike-quoted), \texttt{DeltaSurface} (delta-quoted), and
\texttt{MoneynessSurface} (moneyness-quoted). The configuration is identical to equity,
except the curve lives under \texttt{CommodityVolatility} and requires a
\texttt{PriceCurveId} and \texttt{YieldCurveId}. Commodity forwards are derived from
the \texttt{CommodityPriceCurve} rather than from spot and carry curves: internally, a
\texttt{PriceTermStructureAdapter} converts the price term structure into yield-curve
handles so the existing \texttt{BlackVolatilitySurfaceSvi} machinery is reused directly.

A complete example for a commodity volatility surface using the CorbettaEtAl2019Essvi model variant
on an absolute-strike surface with wildcard expiries and strikes:

\begin{minted}[fontsize=\footnotesize]{xml}
  <CommodityVolatility>
    <CurveId>ICE:B</CurveId>
    <CurveDescription>ICE Brent option surface using ESSVI Robust SVI</CurveDescription>
    <Currency>USD</Currency>
    <VolatilityConfig>
      <StrikeSurface priority="0">
        <QuoteType>ImpliedVolatility</QuoteType>
        <VolatilityType>Lognormal</VolatilityType>
        <Strikes>*</Strikes>
        <Expiries>*</Expiries>
        <TimeInterpolation>CorbettaEtAl2019Essvi</TimeInterpolation>
        <StrikeInterpolation>CorbettaEtAl2019Essvi</StrikeInterpolation>
        <Extrapolation>true</Extrapolation>
        <TimeExtrapolation>Flat</TimeExtrapolation>
        <StrikeExtrapolation>Flat</StrikeExtrapolation>
        <ParametricSmileConfiguration>
          <Parameters>
            <Parameter>
              <Name>k_star</Name>
              <InitialValue>0.0</InitialValue>
              <Calibration>Implied</Calibration>
            </Parameter>
            <Parameter>
              <Name>theta_star</Name>
              <InitialValue>0.0</InitialValue>
              <Calibration>Implied</Calibration>
            </Parameter>
            <Parameter>
              <Name>rho</Name>
              <InitialValue>-0.5</InitialValue>
              <Calibration>Calibrated</Calibration>
            </Parameter>
            <Parameter>
              <Name>psi</Name>
              <InitialValue>1.0</InitialValue>
              <Calibration>Calibrated</Calibration>
            </Parameter>
          </Parameters>
          <Calibration>
            <MaxCalibrationAttempts>100</MaxCalibrationAttempts>
            <ExitEarlyErrorThreshold>0.005</ExitEarlyErrorThreshold>
            <MaxAcceptableError>0.3</MaxAcceptableError>
          </Calibration>
        </ParametricSmileConfiguration>
      </StrikeSurface>
    </VolatilityConfig>
    <DayCounter>A365</DayCounter>
    <Calendar>ICE_FuturesEU</Calendar>
    <FutureConventions>ICE:B</FutureConventions>
    <OptionExpiryRollDays>1</OptionExpiryRollDays>
    <PriceCurveId>Commodity/USD/ICE:B</PriceCurveId>
    <YieldCurveId>Yield/USD/USD-FedFunds</YieldCurveId>
  </CommodityVolatility>
\end{minted}

\subsection{CDS Volatility Surface Example (Mingone, global)}

A complete example for an CDS volatility surface using the Mingone2022EssviGJ model variant:

\begin{minted}[fontsize=\footnotesize]{xml}
  <CDSVolatility>
    <CurveId>RED:2I666VDB8</CurveId>
    <CurveDescription>ITRAXX Europe 5Y index option lognormal vols
      using ESSVI Mingone GJ (global)</CurveDescription>
    <Terms>
      <Term>
        <Label>1Y</Label>
        <Curve>Default/EUR/RED:2I666VDB8_1Y</Curve>
      </Term>
      <Term>
        <Label>3Y</Label>
        <Curve>Default/EUR/RED:2I666VDB8_3Y</Curve>
      </Term>
      <Term>
        <Label>5Y</Label>
        <Curve>Default/EUR/RED:2I666VDB8_5Y</Curve>
      </Term>
      <Term>
        <Label>7Y</Label>
        <Curve>Default/EUR/RED:2I666VDB8_7Y</Curve>
      </Term>
      <Term>
        <Label>10Y</Label>
        <Curve>Default/EUR/RED:2I666VDB8_10Y</Curve>
      </Term>
    </Terms>
    <StrikeSurface priority="0">
      <QuoteType>ImpliedVolatility</QuoteType>
      <VolatilityType>Lognormal</VolatilityType>
      <Strikes>*</Strikes>
      <Expiries>*</Expiries>
      <TimeInterpolation>Mingone2022EssviGJ</TimeInterpolation>
      <StrikeInterpolation>Mingone2022EssviGJ</StrikeInterpolation>
      <Extrapolation>true</Extrapolation>
      <TimeExtrapolation>Flat</TimeExtrapolation>
      <StrikeExtrapolation>Flat</StrikeExtrapolation>
      <ParametricSmileConfiguration>
        <Parameters>
          <Parameter>
            <Name>rho</Name>
            <InitialValue>0.0</InitialValue>
            <Calibration>Calibrated</Calibration>
          </Parameter>
          <Parameter>
            <Name>a</Name>
            <InitialValue>0.0</InitialValue>
            <Calibration>Calibrated</Calibration>
          </Parameter>
          <Parameter>
            <Name>c</Name>
            <InitialValue>0.5</InitialValue>
            <Calibration>Calibrated</Calibration>
          </Parameter>
        </Parameters>
        <Calibration>
          <MaxCalibrationAttempts>50</MaxCalibrationAttempts>
          <ExitEarlyErrorThreshold>0.005</ExitEarlyErrorThreshold>
          <MaxAcceptableError>100.0</MaxAcceptableError>
        </Calibration>
      </ParametricSmileConfiguration>
    </StrikeSurface>
    <DayCounter>A365</DayCounter>
    <Calendar>NullCalendar</Calendar>
    <StrikeType>Spread</StrikeType>
    <QuoteName>2I666VDB8</QuoteName>
    <StrikeFactor>10000</StrikeFactor>
  </CDSVolatility>
\end{minted}

\begin{thebibliography}{10}

\bibitem{Gatheral_2004} Gatheral, J., A parsimonious arbitrage-free implied volatility parameterization with application to the valuation of volatility derivatives, Presentation at Global Derivatives \& Risk Management, 2004

\bibitem{Gatheral_2012} Gatheral, J., Jacquier A., Arbitrage-free SVI volatility surfaces, Quantitative Finance, 2012

\bibitem{Hendriks_2017} Hendriks S., Martini C., The extended SSVI volatility surface, 2017

\bibitem{Corbetta_2019} Corbetta J., Cohort P., Laachir I., Martini C., Robust calibration and arbitrage-free interpolation of SSVI slices, 2019

\bibitem{Mingone_2022} Mingone A., No arbitrage global parametrization for the eSSVI volatility surface, 2022

\end{thebibliography}


\end{document}
